% ══════════════════════════════════════════════════════════════
% Chapter 4 — Market Microstructure and Dealer Positioning
% ══════════════════════════════════════════════════════════════

\chapter{Market Microstructure and Dealer Positioning}
\label{ch:microstructure}

\begin{projectconnection}[Project Connection --- Phase 2 \& Phase 4B]
This chapter teaches the mechanics of dealer positioning that create
your features.  VaR utilization and factor concentration
(Phase~2, Project~1) reflect dealer constraints from
Chapter~\ref{ch:intermediary}'s theory.  Dealer Greeks---gamma, vanna,
charm (Phase~4B, Project~2)---create mechanical hedging flows that
move prices.  Understanding \emph{why} these flows exist is the
difference between engineering meaningful features and fitting noise.
\end{projectconnection}


% ──────────────────────────────────────────────────────────────
\section{Dealer Balance Sheets and Inventory}
\label{sec:dealer-inventory}
% ──────────────────────────────────────────────────────────────

\begin{prereq}[Long and Short Positions]
A \textbf{long} position means you own the asset and profit when its
price rises.  A \textbf{short} position means you have sold an asset
you do not own (by borrowing it) and profit when the price falls.
A dealer who sells a call option to a client is \emph{short} the call;
a dealer who buys a bond from a client is \emph{long} the bond.
\end{prereq}

\begin{prereq}[Bid-Ask Spread]
The \textbf{bid} is the highest price a dealer will pay to buy an
asset; the \textbf{ask} (or offer) is the lowest price a dealer will
accept to sell.  The \textbf{bid-ask spread} $= \text{ask} - \text{bid}$
is the dealer's compensation for providing immediacy.  Typical spreads:
$\sim$1 cent for liquid large-cap equities, $\sim$1--5 basis points
for on-the-run Treasuries, $\sim$10--50 basis points for corporate
bonds.
\end{prereq}

Market-makers exist because buyers and sellers rarely arrive at the same
time.  A client wants to sell \$50M of investment-grade bonds today;
the natural buyer may not appear for days.  The dealer steps in, buys
the bonds at the bid price, warehouses them, and hopes to sell at the
ask before the market moves against the position.

This creates \textbf{inventory}---risk the dealer holds involuntarily
as a byproduct of providing liquidity.  Inventory has two costs:

\begin{enumerate}[nosep]
  \item \textbf{Carry cost:} The position uses capital and consumes
        balance sheet.  The dealer must fund the position (repo,
        margin) and hold regulatory capital against it.
  \item \textbf{Market risk:} The position contributes to the desk's
        VaR (Chapter~\ref{ch:risk-systems}).  Large inventory pushes
        VaR utilization toward the limit, constraining the desk's
        ability to make further markets.
\end{enumerate}

\begin{keyidea}[Inventory Links to VaR Utilization]
A desk with high inventory has high VaR.  High VaR relative to the
limit means high VaR utilization.  High VaR utilization constrains
future market-making, reducing liquidity provision.  This is the
micro-level mechanism behind the procyclical leverage documented by
Adrian and Shin (2014) and the intermediary constraints priced by
He, Krishnamurthy (2013)---concepts from Chapter~\ref{ch:intermediary}.

\textbf{For your project:} When VaR utilization rises rapidly
(Section~\ref{sec:var-utilization} of Chapter~\ref{ch:risk-systems}),
it signals that dealers are accumulating involuntary inventory.
The subsequent forced reduction of that inventory creates
predictable price pressure.
\end{keyidea}


% ──────────────────────────────────────────────────────────────
\section{Delta Hedging}
\label{sec:delta-hedging}
% ──────────────────────────────────────────────────────────────

\begin{prereq}[What Is an Option?]
An \textbf{option} gives the holder the right---but not the
obligation---to buy (call) or sell (put) an underlying asset at a
specified \textbf{strike price} $K$ on or before an
\textbf{expiration date} $T$.

\textbf{Payoff at expiration:}
\begin{itemize}[nosep]
  \item Call: $\max(S_T - K,\; 0)$.  Profits when $S_T > K$.
  \item Put:  $\max(K - S_T,\; 0)$.  Profits when $S_T < K$.
\end{itemize}
Before expiration, an option's value $V(S,t)$ depends on the current
price $S$, time to expiration $T-t$, volatility $\sigma$, interest
rate $r$, and the strike $K$.  The Black-Scholes model provides the
canonical pricing formula, but the key insight for this chapter is
simpler: \emph{option value changes when the underlying moves.}
\end{prereq}

Options dealers do not speculate on direction---they earn the bid-ask
spread on options and hedge away the directional risk.  The primary
hedge is \textbf{delta hedging}.

\begin{definition}[Delta]
The \textbf{delta} of an option is the first derivative of its value
with respect to the underlying price:
\begin{equation}
\label{eq:delta-def}
\Delta \;=\; \frac{\partial V}{\partial S}
\end{equation}
\begin{itemize}[nosep]
  \item $V$ --- option value.
  \item $S$ --- price of the underlying asset.
  \item Call delta: $\Delta \in [0,\,1]$.  A deep in-the-money call
        has $\Delta \approx 1$; a far out-of-the-money call has
        $\Delta \approx 0$.
  \item Put delta: $\Delta \in [-1,\,0]$.  Same logic, opposite sign.
  \item $\Delta$ changes as $S$ moves---it is not constant.
\end{itemize}
\end{definition}

\textbf{Delta hedging in practice:} If a dealer sells a call option
with $\Delta = 0.50$ on 100 shares, the dealer is short delta by
$0.50 \times 100 = 50$ shares.  To hedge, the dealer buys 50 shares
of the underlying.  Now, if the stock rises by \$1, the option loss
($-\$50$) is offset by the stock gain ($+\$50$).

The problem: delta changes as $S$ moves.  If the stock rises, the
call's delta increases (say to 0.60), and the dealer must buy
10 more shares.  If the stock falls, the call's delta decreases
(say to 0.40), and the dealer must sell 10 shares.

\begin{keyidea}[Delta Hedging Creates Mechanical Flows]
Delta hedging generates \textbf{predictable, mechanical trading} in
the underlying asset.  When the stock moves, dealers must rebalance.
The direction of rebalancing depends on whether dealers are net long
or net short \emph{gamma}---the topic of the next section.

These flows are not information-driven.  They are constraint-driven:
the dealer \emph{must} hedge regardless of any view on future
direction.  This is precisely the kind of non-fundamental flow that
can create predictable short-term return patterns.
\end{keyidea}

\begin{workedexample}{Delta Hedging a Short Call Position}
A dealer sells 1{,}000 at-the-money call contracts on a stock at
$S = \$100$, $K = \$100$.  Each contract covers 100 shares.
Delta of each call: $\Delta = 0.50$.

\textbf{Initial hedge:}
Total delta exposure $= -1{,}000 \times 100 \times 0.50
  = -50{,}000$ shares (short delta).
The dealer buys 50{,}000 shares to hedge.

\textbf{Stock rises to \$105:}
New delta $\approx 0.62$.
Required hedge $= 1{,}000 \times 100 \times 0.62 = 62{,}000$ shares.
The dealer must buy $62{,}000 - 50{,}000 = 12{,}000$ additional shares.

\textbf{Stock falls to \$95:}
New delta $\approx 0.38$.
Required hedge $= 1{,}000 \times 100 \times 0.38 = 38{,}000$ shares.
The dealer must sell $50{,}000 - 38{,}000 = 12{,}000$ shares.

In both cases, the dealer trades \emph{in the direction of the move}:
buying on the way up, selling on the way down.  This amplifies the
price change---the hallmark of a short-gamma position.
\end{workedexample}


% ──────────────────────────────────────────────────────────────
\section{Gamma and the Gamma-Flip Level}
\label{sec:gamma}
% ──────────────────────────────────────────────────────────────

\begin{definition}[Gamma]
The \textbf{gamma} of an option is the second derivative of its
value with respect to the underlying price, or equivalently, the
rate of change of delta:
\begin{equation}
\label{eq:gamma-def}
\Gamma \;=\; \frac{\partial^2 V}{\partial S^2}
  \;=\; \frac{\partial \Delta}{\partial S}
\end{equation}
\begin{itemize}[nosep]
  \item $\Gamma > 0$ always for plain vanilla options
        (both calls and puts).
  \item Gamma is highest for at-the-money options near expiration.
  \item A dealer who \emph{buys} options is \textbf{long gamma}
        ($\Gamma > 0$).
  \item A dealer who \emph{sells} options is \textbf{short gamma}
        ($\Gamma < 0$ for the dealer's portfolio).
\end{itemize}
\end{definition}

The sign of aggregate dealer gamma determines whether hedging flows
stabilize or destabilize prices.

\begin{intuition}[Shock Absorbers vs.\ Shock Amplifiers]
\textbf{Long-gamma dealers are shock absorbers.}  When the underlying
rises, their delta increases, so they \emph{sell} into the rally
(locking in gains).  When it falls, their delta decreases, so they
\emph{buy} the dip.  Their hedging dampens price moves.

\textbf{Short-gamma dealers are shock amplifiers.}  When the underlying
rises, their negative gamma means they must \emph{buy more} to stay
hedged.  When it falls, they must \emph{sell more}.  Their hedging
amplifies price moves.

The effect is proportional to the \emph{aggregate} net gamma across
all dealers.  One dealer's long gamma can cancel another's short gamma.
What matters for price impact is the signed sum.
\end{intuition}

\vspace{0.5em}

% ── TikZ Diagram: Long vs Short Gamma Hedging ──
\begin{figure}[ht]
\centering
\begin{tikzpicture}[
  node distance=0.5cm and 2.5cm,
  every node/.style={font=\small},
  flowblock/.style={draw, rounded corners, minimum height=0.8cm,
    minimum width=2.8cm, align=center, font=\small},
  arrow/.style={-{Stealth[length=2.5mm]}, thick},
]

% ── Left column: Long Gamma ──
\node[font=\bfseries\small, text=intgreen] (ltitle) {Net Long Gamma};
\node[flowblock, fill=green!8, below=0.4cm of ltitle] (lshock)
  {Price rises by $\delta$};
\node[flowblock, fill=green!8, below=0.4cm of lshock] (ldelta)
  {$\Delta$ increases\\(calls more in-the-money)};
\node[flowblock, fill=green!8, below=0.4cm of ldelta] (lhedge)
  {Dealer \textbf{sells} shares\\to reduce hedge};
\node[flowblock, fill=green!15, below=0.4cm of lhedge] (lresult)
  {Selling pressure\\dampens rally};

\draw[arrow, intgreen] (lshock) -- (ldelta);
\draw[arrow, intgreen] (ldelta) -- (lhedge);
\draw[arrow, intgreen] (lhedge) -- (lresult);

% ── Right column: Short Gamma ──
\node[font=\bfseries\small, text=warnred, right=4cm of ltitle]
  (rtitle) {Net Short Gamma};
\node[flowblock, fill=red!8, below=0.4cm of rtitle] (rshock)
  {Price rises by $\delta$};
\node[flowblock, fill=red!8, below=0.4cm of rshock] (rdelta)
  {$\Delta$ increases\\(short calls deeper ITM)};
\node[flowblock, fill=red!8, below=0.4cm of rdelta] (rhedge)
  {Dealer \textbf{buys} shares\\to increase hedge};
\node[flowblock, fill=red!15, below=0.4cm of rhedge] (rresult)
  {Buying pressure\\amplifies rally};

\draw[arrow, warnred] (rshock) -- (rdelta);
\draw[arrow, warnred] (rdelta) -- (rhedge);
\draw[arrow, warnred] (rhedge) -- (rresult);

% ── Center label ──
\node[font=\footnotesize\itshape, text=gray,
  below=0.3cm of $(lresult)!0.5!(rresult)$]
  {Symmetric logic applies when prices fall};

\end{tikzpicture}
\caption{Dealer hedging flows under net long gamma (stabilizing)
vs.\ net short gamma (destabilizing).  The mechanism is symmetric
for price declines: long-gamma dealers buy dips; short-gamma dealers
sell into declines.}
\label{fig:gamma-flows}
\end{figure}

\subsection{Aggregate Gamma and the Gamma-Flip Level}

Individual option positions have positive gamma.  But \emph{dealers}
can be net long or net short gamma depending on whether they are net
buyers or net sellers of options.  The \textbf{aggregate net dealer
gamma} is the dollar-weighted sum of gamma across all dealer
positions:

\begin{equation}
\label{eq:agg-gamma}
\Gamma_{\text{agg}} \;=\; \sum_{i} n_i \cdot \Gamma_i \cdot S^2
  \cdot 0.01
\end{equation}
\begin{itemize}[nosep]
  \item $n_i$ --- net dealer position in contract $i$ (positive if
        dealer is long, negative if short).
  \item $\Gamma_i$ --- gamma of contract $i$.
  \item $S$ --- underlying price.
  \item The $S^2 \times 0.01$ normalization converts to ``dollars of
        stock a dealer must trade per 1\% move in $S$.''
\end{itemize}

The \textbf{gamma-flip level} (also called the \textbf{gamma-neutral
level}) is the price $S^*$ at which $\Gamma_{\text{agg}} = 0$.  Below
$S^*$, dealers are typically net short gamma (puts gain gamma as the
underlying falls); above $S^*$, dealers are typically net long gamma.

\begin{keyidea}[The Gamma-Flip Level as a Regime Boundary]
When the underlying is above the gamma-flip level, dealers' hedging
dampens volatility (long gamma).  When it crosses below the flip
level, the regime switches: dealers' hedging amplifies volatility
(short gamma).  This transition---from shock absorber to shock
amplifier---often coincides with sharp increases in realized
volatility and intraday price swings.

\textbf{For your project (Phase~4B):} If you compute aggregate
dealer gamma from SecDB's book-level Greeks, the gamma-flip level
becomes a natural feature.  Distance to the flip level, and whether
the market recently crossed it, can predict volatility regime changes.
\end{keyidea}


% ──────────────────────────────────────────────────────────────
\section{Higher-Order Greeks}
\label{sec:higher-greeks}
% ──────────────────────────────────────────────────────────────

Gamma is the primary driver of mechanical hedging flows, but three
higher-order sensitivities create additional predictable trading
patterns.

\begin{prereq}[Partial Derivatives]
The Greeks are partial derivatives of option value $V(S, \sigma, t)$
with respect to different variables.  Delta is $\partial V / \partial S$.
Gamma is $\partial^2 V / \partial S^2$.  The higher-order Greeks below
are cross-partials or time derivatives.  You do not need to derive
them---just understand what they measure.
\end{prereq}

\begin{definition}[Vega, Vanna, and Charm]
\textbf{Vega} --- sensitivity of option value to implied volatility:
\begin{equation}
\label{eq:vega-def}
\mathcal{V} \;=\; \frac{\partial V}{\partial \sigma}
\end{equation}
Vega is always positive for long option positions.  When implied
volatility rises by 1 percentage point, the option value changes by
approximately $\mathcal{V}$.

\textbf{Vanna} --- cross-sensitivity of delta to volatility (or
equivalently, sensitivity of vega to the underlying price):
\begin{equation}
\label{eq:vanna-def}
\text{Vanna} \;=\; \frac{\partial^2 V}{\partial S\, \partial \sigma}
  \;=\; \frac{\partial \Delta}{\partial \sigma}
  \;=\; \frac{\partial \mathcal{V}}{\partial S}
\end{equation}
Vanna matters on days when both $S$ and $\sigma$ move together---i.e.,
during market stress, when prices drop and volatility spikes
simultaneously.

\textbf{Charm} (delta decay) --- sensitivity of delta to the passage
of time:
\begin{equation}
\label{eq:charm-def}
\text{Charm} \;=\; -\frac{\partial \Delta}{\partial t}
  \;=\; -\frac{\partial^2 V}{\partial S\, \partial t}
\end{equation}
Charm captures how delta changes overnight, even if $S$ and $\sigma$
are unchanged.  As options approach expiration, at-the-money delta
moves toward 0.50, and out-of-the-money delta decays toward 0.
\begin{itemize}[nosep]
  \item $\mathcal{V}$ drives hedging when volatility moves
        (earnings, FOMC).
  \item Vanna drives hedging when price and volatility move
        \emph{simultaneously} (sell-offs with vol spikes).
  \item Charm drives end-of-day hedging as time passes---even on
        quiet days.
\end{itemize}
\end{definition}

\begin{keyidea}[Which Greek Matters When]
\begin{center}
\small
\begin{tabular}{@{}llp{6cm}@{}}
\toprule
\textbf{Greek} & \textbf{Trigger} & \textbf{Hedging Flow} \\
\midrule
Gamma & Price moves &
  Dominant hedging flow.  Direction depends on dealer net gamma sign. \\
Vanna & Vol move + price move &
  Large on event days (earnings, FOMC) when $S$ and $\sigma$ move
  together.  Creates asymmetric flows in sell-offs. \\
Charm & Time passing &
  Predictable end-of-day flow as dealers adjust for overnight
  delta decay.  Strongest near expiration. \\
Vega & Vol moves alone &
  Matters for vol surface trading but less relevant for
  directional flow in the underlying. \\
\bottomrule
\end{tabular}
\end{center}

For Phase~4B, gamma is the primary feature.  Add vanna as a
feature on event days (scheduled vol events like FOMC, NFP, CPI)
and charm for modeling last-hour flows.
\end{keyidea}

\begin{workedexample}{Vanna Flow During a Sell-Off}
Suppose dealers are short a large book of put options on the S\&P~500.
The market drops 2\% and the VIX spikes 5 vol points.

\textbf{Gamma effect:} The puts' delta becomes more negative
(closer to $-1$).  Dealers who are short these puts have positive
delta exposure increasing---they must \emph{sell} index futures to
re-hedge.  This amplifies the decline (short gamma).

\textbf{Vanna effect:} The simultaneous volatility spike further
increases the magnitude of put delta (higher vol $\Rightarrow$
out-of-the-money puts gain delta).  Dealers must sell \emph{even more}
futures than gamma alone would dictate.

During the COVID crash of March 2020, the compounding of gamma and
vanna hedging flows was widely cited as an amplifier of the sell-off.
The vanna effect is often larger than the gamma effect during
high-volatility episodes because $\sigma$ and $S$ move
simultaneously.
\end{workedexample}


% ──────────────────────────────────────────────────────────────
\section{Fire Sales: Coval and Stafford (2007)}
\label{sec:fire-sales}
% ──────────────────────────────────────────────────────────────

The theoretical framework from Chapter~\ref{ch:intermediary} predicts
that constrained intermediaries create predictable price pressure.
Coval and Stafford (2007) provide the cleanest empirical evidence.

\begin{keyresult}[Fire Sales and Reversals (Coval and Stafford, 2007, JFE)]
Using mutual fund holdings data from 1980--2003, Coval and Stafford
identify stocks subject to flow-driven forced selling (funds
experiencing large outflows that mechanically liquidate existing
positions).

\textbf{Key findings:}
\begin{itemize}[nosep]
  \item Stocks in the heaviest forced-selling decile experience
        cumulative abnormal returns of approximately $-10\%$ during
        the selling quarter ($t$-stat $\approx -3.7$).
  \item Reversals emerge gradually: $+2\%$ in the following quarter,
        approximately $+8\%$ over the subsequent year, with full mean
        reversion at roughly 18--20 months.
  \item A long-short strategy trading against anticipated forced
        sales earns abnormal returns exceeding $10\%$ annually.
  \item Funds in the bottom outflow decile are roughly twice as
        likely to liquidate positions as funds with normal flows.
\end{itemize}
The mechanism: capital outflows force funds to sell regardless of
fundamental value, pushing prices below fair value.  Prices revert
as the forced selling subsides and other investors provide liquidity.
\end{keyresult}

The Coval-Stafford mechanism applies directly to dealer banks, but
through a different constraint.  Mutual funds are forced to sell by
\emph{redemptions}.  Dealers are forced to sell by \emph{VaR limits}.

\begin{keyidea}[VaR Limits as the Dealer Fire-Sale Trigger]
In Chapter~\ref{ch:risk-systems}, Section~\ref{sec:var-utilization},
you learned that VaR utilization approaching 100\% forces desks to
reduce positions.  This is the dealer-bank analog of Coval-Stafford
fund outflows:

\begin{center}
\small
\begin{tabular}{@{}lll@{}}
\toprule
 & \textbf{Mutual Fund} & \textbf{Dealer Bank} \\
\midrule
Constraint & Redemptions (outflows) & VaR limit \\
Trigger & Large capital outflows & VaR utilization $> 90\%$ \\
Forced action & Liquidate existing positions & Reduce positions
  or hedge \\
Price impact & $\approx -10\%$ CAR & Expected similar magnitude \\
Reversal & 18--20 months to mean revert & Your project tests this \\
\bottomrule
\end{tabular}
\end{center}

The key difference: Coval-Stafford used public fund holdings
(13F filings, quarterly).  You have daily VaR utilization from
SecDB---a higher-frequency, more granular version of the same
constraint.
\end{keyidea}


% ──────────────────────────────────────────────────────────────
\section{Dealer Gamma and Intraday Momentum}
\label{sec:intraday-momentum}
% ──────────────────────────────────────────────────────────────

Baltussen, Da, Lammers, and Martens (2021, JFE) connect the
gamma-hedging mechanics from Section~\ref{sec:gamma} to a
specific, tradeable pattern: \textbf{market intraday momentum}.

\begin{keyresult}[Intraday Momentum from Gamma Hedging
  (Baltussen et al., 2021, JFE)]
Using intraday returns on 60+ futures contracts across equities,
bonds, commodities, and currencies from 1974--2020:
\begin{itemize}[nosep]
  \item The return in the \textbf{last 30 minutes} before market
        close is positively predicted by the return during the
        rest of the day (open to 30 minutes before close).
  \item The predictability is economically large: a simple trading
        strategy (go long if the day's return is positive, short if
        negative, held for the last 30 minutes) produces annualized
        Sharpe ratios between \textbf{0.87 and 1.73} at the
        asset-class level.
  \item The effect is strongest on days with large absolute returns
        and reverts over the following days.
  \item The authors link the pattern to gamma hedging demand from
        options market-makers and leveraged ETF rebalancing: both
        must trade in the direction of the day's move near the close.
\end{itemize}
\end{keyresult}

\textbf{The mechanism in detail:}  Suppose the S\&P~500 is up 1.5\%
by 3:30pm.  Dealers who are net short gamma on index options must buy
more futures to rebalance their delta hedge (Section~\ref{sec:gamma}).
Leveraged ETFs that target $2\times$ or $3\times$ daily returns must
also buy to restore their target leverage ratio.  Both flows
concentrate near the close, pushing the index further in the
direction of the day's move.

\begin{keyidea}[Why Proprietary Data Beats Public GEX]
Public ``Gamma Exposure'' (GEX) estimates---available from services
like SqueezeMetrics and SpotGamma---reconstruct aggregate dealer gamma
from exchange-reported options data.  These estimates have three
problems:

\begin{enumerate}[nosep]
  \item \textbf{Missing OTC volume:} Public data covers only
        exchange-listed options.  A substantial portion of dealer
        gamma sits in over-the-counter contracts that are invisible
        to public data.
  \item \textbf{Unknown dealer side:} Public data shows open interest
        and volume but not \emph{who} is long and who is short.
        Assigning dealer vs.\ customer sides requires assumptions
        that introduce sign errors.
  \item \textbf{Incomplete expiration coverage:} Different providers
        include different subsets of expirations (0DTE, weeklies,
        monthlies), producing conflicting estimates for the same
        underlying on the same day.
\end{enumerate}

SecDB provides the dealer's \emph{actual} book-level Greeks with
correct sign---no assumptions about positioning needed.  This is
the core advantage of Project~2 (Phase~4B).
\end{keyidea}


% ──────────────────────────────────────────────────────────────
\section{Option Expiration Effects}
\label{sec:expiration-effects}
% ──────────────────────────────────────────────────────────────

Dealer gamma effects are not constant---they intensify near option
expiration dates.  Ni, Pearson, and Poteshman (2005, JFE) document
the most direct evidence.

\begin{keyresult}[Stock Pinning at Strikes (Ni, Pearson, and
  Poteshman, 2005, JFE)]
On option expiration dates, closing prices of stocks with listed
options \textbf{cluster at strike prices}.  On each expiration date,
the returns of optionable stocks are altered by an average of at
least \textbf{16.5 basis points}, translating into aggregate
market-capitalization shifts on the order of \textbf{\$9 billion}
(in their 1996--2002 sample).

The mechanism: delta-hedging market-makers with net long option
positions sell stock above the nearest strike and buy below it,
driving the closing price toward the strike.  The effect reverses
for net short positions, creating ``declustering.''
\end{keyresult}

The expiration-pinning effect is relevant to your project for two
reasons.  First, it is direct evidence that dealer hedging moves
prices in a predictable, non-informational way.  Second, the rise of
zero-days-to-expiration (0DTE) options---which now account for over
40\% of S\&P~500 options volume---means that this pinning effect
occurs \emph{every trading day}, not just on monthly expiration
Fridays.

\begin{workedexample}{Gamma Near Expiration}
Consider an at-the-money call with $K = \$100$, $S = \$100$, and
1 day to expiration.  Its gamma is extremely high: a small move in
$S$ causes a large change in delta.

If $S$ rises to \$101 (1\% move), the call's delta jumps from
$\approx 0.50$ to $\approx 0.85$.  A dealer short this call must
buy 35 additional shares per contract---a much larger rebalance than
for an option with 30 days to expiration, where the same move might
shift delta by only 5 shares.

Near expiration, gamma is concentrated at the strike.  This is why
pinning effects are strongest at strikes with the most open interest,
and why dealers' hedging flows peak in the final hours of expiration
day.
\end{workedexample}


% ──────────────────────────────────────────────────────────────
\section{Caveats and Pitfalls}
\label{sec:microstructure-caveats}
% ──────────────────────────────────────────────────────────────

The dealer-positioning channel is real, but three caveats
from the literature require careful handling.

\subsection{Implied Volatility Spread as a Borrow-Fee Proxy}

\begin{keyresult}[IVS Reflects Borrow Fees, Not Information
  (Muravyev, Pearson, and Pollet, 2025, JFE)]
The implied-volatility spread (call IV minus put IV) and implied
volatility skew are widely used as signals for future stock returns.
Muravyev, Pearson, and Pollet show that this apparent predictability
is largely a proxy for \textbf{stock borrow fees}:
\begin{itemize}[nosep]
  \item The return predictability from IV spread and skew decreases
        by at least \textbf{two-thirds} when high-borrow-fee stocks
        are excluded from the sample.
  \item The mechanism: standard implied volatility calculations
        assume zero borrowing costs.  When borrow fees are high
        (hard-to-borrow stocks), this assumption inflates put IV
        relative to call IV, creating a spurious ``signal.''
  \item Strategies that appear profitable in backtests become
        unprofitable once actual borrowing costs are included.
\end{itemize}
\end{keyresult}

\begin{warning}[Control for Short Interest When Using Options Data]
If your Project~2 features touch equities (even through index
components), any signal derived from implied volatility differences
or options order flow may be contaminated by borrow fees.
Muravyev, Pearson, and Pollet (2025) show that at least two-thirds
of the return predictability from IV spread and skew disappears when
high-borrow-fee stocks are excluded.

\textbf{Mitigation:}
\begin{enumerate}[nosep]
  \item If available, include stock borrow fees as a control variable
        in your regressions.
  \item At minimum, add short interest as a percentage of float as a
        control.
  \item Test whether your signal survives when restricted to the
        subset of stocks with low short interest and low borrow fees.
\end{enumerate}
This caveat is less binding for rates and FX futures (Project~1's
primary asset classes), where borrow fees are negligible.  It becomes
critical if Project~2 extends to single-stock options.
\end{warning}

\subsection{Gamma Fragility: Sign Flips Intraday}

Barbon and Buraschi (2021) provide granular evidence on how dealer
gamma affects returns at the stock level, using options data from
Nasdaq exchanges (ISE, GEMX, PHLX) over 2010--2020.

\begin{keyresult}[Gamma Fragility (Barbon and Buraschi, 2021)]
A one-standard-deviation increase in the dealer gamma imbalance is
associated with:
\begin{itemize}[nosep]
  \item A decrease of \textbf{5--25 basis points} in absolute daily
        returns for single stocks (stabilization from long gamma).
  \item A decrease of more than \textbf{20 basis points}
        ($\approx$20\% of a standard deviation) in absolute daily
        returns for indexes.
  \item A decrease in intraday return autocorrelation of
        \textbf{1.5--2\%} per standard deviation of gamma imbalance.
  \item The effect is significantly larger for \textbf{illiquid}
        stocks, where dealer hedging flows represent a larger share
        of total volume.
\end{itemize}
The fragility arises because gamma imbalances can flip sign quickly
as new option trades execute, existing options expire, and the
underlying price crosses strike boundaries.
\end{keyresult}

\begin{warning}[Dealer Gamma Changes Intraday]
The gamma estimates you compute from SecDB's morning risk run may be
stale by the afternoon close.  Three sources of intraday change:
\begin{enumerate}[nosep]
  \item \textbf{New option trades} change dealer positions throughout
        the day.  A single large customer order can flip the desk's
        net gamma sign.
  \item \textbf{Price movement} shifts which options are at-the-money
        (where gamma is highest).  A 2\% intraday move can
        redistribute gamma across strikes.
  \item \textbf{Time decay} (charm) alters delta and gamma
        continuously, especially for near-expiration options.
\end{enumerate}

\textbf{Mitigation for Phase~4B:}
\begin{itemize}[nosep]
  \item Use end-of-day Greeks (after the close), not morning
        estimates, when constructing daily features.
  \item Treat the gamma feature as a \emph{regime indicator}
        (long gamma vs.\ short gamma) rather than a precise dollar
        estimate.
  \item Include the absolute magnitude of gamma as a separate feature
        measuring the \emph{intensity} of hedging flows, regardless
        of direction.
\end{itemize}
\end{warning}


% ──────────────────────────────────────────────────────────────
\section{Summary}
\label{sec:micro-summary}
% ──────────────────────────────────────────────────────────────

\begin{table}[h]
\centering
\small
\begin{tabular}{@{}p{3cm}p{4.5cm}p{6cm}@{}}
\toprule
\textbf{Concept} & \textbf{Mechanism} & \textbf{Project Relevance} \\
\midrule
Dealer inventory &
  Market-making creates involuntary positions that consume VaR &
  Links to VaR utilization features
  (Ch.~\ref{ch:risk-systems}) \\[6pt]

Delta hedging &
  Dealers trade the underlying to offset option exposure;
  $\Delta = \partial V / \partial S$ &
  Creates the mechanical flows underlying all Greek-based
  features \\[6pt]

Gamma &
  $\Gamma = \partial^2 V / \partial S^2$; determines whether
  hedging stabilizes or amplifies prices &
  Primary feature for Project~2 (Phase~4B); gamma-flip level
  as regime boundary \\[6pt]

Vanna, charm &
  Cross-partials that drive hedging on event days (vanna) and
  near the close (charm) &
  Secondary features for Phase~4B; add on scheduled-event
  days \\[6pt]

Fire sales (Coval-Stafford) &
  Forced selling $\Rightarrow$ $\approx -10\%$ CAR
  $\Rightarrow$ reversal over 18--20 months &
  VaR limit breaches are the dealer analog;
  Project~1 tests this with daily data \\[6pt]

Intraday momentum (Baltussen et al.) &
  Gamma hedging predicts last-30-min returns;
  SR $\in [0.87, 1.73]$ &
  Direct test for Project~2 using SecDB book Greeks \\[6pt]

Gamma fragility (Barbon-Buraschi) &
  Gamma imbalance dampens/amplifies returns by 5--25 bps
  per $\sigma$; stronger for illiquid assets &
  Use gamma as a regime indicator, not a precise estimate;
  effect is fragile intraday \\[6pt]

IVS caveat (Muravyev et al.) &
  IV spread/skew proxies borrow fees; $\ge 2/3$ of
  predictability is spurious &
  Control for short interest if touching equity options \\
\bottomrule
\end{tabular}
\caption{Chapter~\ref{ch:microstructure} key concepts and project
relevance.}
\label{tab:ch4-summary}
\end{table}

\vspace{1em}

\noindent\textbf{Key takeaways:}
\begin{enumerate}[nosep]
  \item Dealers accumulate inventory by providing liquidity.
        Inventory consumes VaR, linking microstructure to the
        risk-system outputs in Chapter~\ref{ch:risk-systems}.
  \item Delta hedging creates mechanical, non-informational
        trading flows.  The direction depends on the sign of
        aggregate dealer gamma.
  \item Net long gamma stabilizes prices (sell rallies, buy dips).
        Net short gamma destabilizes them (buy rallies, sell dips).
        The gamma-flip level is a regime boundary.
  \item Coval and Stafford (2007) show forced selling creates
        $\approx -10\%$ abnormal returns with subsequent reversal.
        VaR limits are the dealer-bank version of the same constraint.
  \item Baltussen et al.\ (2021) show gamma hedging predicts
        last-30-minute returns with Sharpe ratios of 0.87--1.73.
        SecDB book Greeks give you the correct sign that public
        data cannot.
  \item Implied-volatility signals may proxy borrow fees
        (Muravyev et al., 2025).  Control for short interest when
        using options-derived features on equities.
  \item Dealer gamma is fragile: it changes intraday and flips
        sign quickly.  Treat it as a regime indicator, not a precise
        estimate.
\end{enumerate}

\noindent The next chapter (Chapter~\ref{ch:regularized}) introduces
regularized linear models---the ridge and lasso baselines that every
ML signal must beat before you trust it.
